%%
%% ASSETS Poster Extended Abstract
%% PP Application: Task Management for ADHD via Cognitive Strengths
%%
\documentclass[sigconf]{acmart}

\AtBeginDocument{%
  \providecommand\BibTeX{{%
  Bib\TeX}}}

\setcopyright{acmlicensed}
\copyrightyear{2025}
\acmYear{2025}
\acmDOI{XXXXXXX.XXXXXXX}
\acmConference[ASSETS '25]{The 27th International ACM SIGACCESS Conference on
  Computers and Accessibility}{October 26--29, 2025}{Denver, CO, USA}
\acmISBN{978-1-4503-XXXX-X/2025/10}

\begin{document}

\title[PP: ADHD Task Management via Cognitive Strengths]{PP Application: A Task
  Management Tool for People with ADHD That Amplifies Cognitive Advantages to
  Correct Disadvantages}

\author{Jingyu Lei}
\affiliation{%
  \institution{University of Cambridge}
  \department{Department of Architecture}
  \city{Cambridge}
  \country{United Kingdom}}
\email{jl2557@ca.ac.uk}

\author{Lan Xiao}
\affiliation{%
  \institution{University College London}
  \department{Department of Computer Science, DDI}
  \city{London}
  \country{United Kingdom}}
\email{l.xiao.22@ucl.ac.uk}

\renewcommand{\shortauthors}{Lei}

\begin{abstract}
Neurodiversity refers to how people's brains work and interpret information, and
ADHD is a common form of neurodiversity. Characterised by age-inappropriate
hyperactivity, impulsivity, and inattention, it is widely defined as a cognitive
deficit. However, little scientific literature has focused on its positive
effects. Current digital applications that assist people with ADHD in managing
tasks are mostly based on control tools---reminders, timers, and behavioural
correction---which reinforce the negative mental models of people with ADHD and
show limited long-term effectiveness. Previous research has shown that adults
with ADHD have a strong capacity for unconscious thought and six identifiable
cognitive strengths. This paper presents the \textbf{PP (Positive and Power)
application}, a task management tool that uses each of those six cognitive
strengths to circumvent its corresponding executive impairment, because
``people with differences may not need a cure; they need help and accommodation.''
\end{abstract}

\begin{CCSXML}
<ccs2012>
 <concept>
  <concept_id>10003120.10003121</concept_id>
  <concept_desc>Human-centered computing~Human computer interaction (HCI)</concept_desc>
  <concept_significance>500</concept_significance>
 </concept>
 <concept>
  <concept_id>10003120.10003121.10003124</concept_id>
  <concept_desc>Human-centered computing~Interactive systems and tools</concept_desc>
  <concept_significance>300</concept_significance>
 </concept>
 <concept>
  <concept_id>10003120.10003121.10003125.10010595</concept_id>
  <concept_desc>Human-centered computing~User interface programming</concept_desc>
  <concept_significance>100</concept_significance>
 </concept>
</ccs2012>
\end{CCSXML}

\ccsdesc[500]{Human-centered computing~Human computer interaction (HCI)}
\ccsdesc[300]{Human-centered computing~Interactive systems and tools}
\ccsdesc[100]{Human-centered computing~User interface programming}

\keywords{ADHD, Neurodiversity, Cognitive advantages, Digital application,
  Assistive technology, Task management}

\maketitle

\section{Introduction}

Behavioural characteristics of ADHD do not exist in a binary form of normal and
abnormal~\cite{sedgwick2019}. A person with ADHD experiences deficits in
cognitive control, sensitivity to rewards, and time perception, while
simultaneously displaying high energy, creativity, hyperfocus, empathy, and
willingness to help others. The core deficit is \emph{executive dysfunction}
\cite{nigg2005}---a broad range of top-down cognitive processes enabling
flexible, goal-directed behaviour---which significantly affects the daily social
activities of adults with ADHD.

Current psychosocial interventions include digital apps that assist with task
execution. A systematic review of 355 app-related studies found that the
majority of effectiveness outcomes were negative~\cite{pasarelu2020}: benefits
did not persist after the intervention period, and users reported increased
anxiety when separated from the tool. More critically, these control-oriented
tools reinforce the negative mental model that ``I need to be managed,''
undermining self-esteem.

Brown's model of ADHD executive functioning identifies six impairment
dimensions---activation, concentration, effort, emotion, memory, and
execution. A complementary qualitative study defines six positive core themes in
ADHD: \textit{cognitive vitality}, \textit{courage}, \textit{energy},
\textit{humanity}, \textit{resilience}, and \textit{transcendence}
\cite{schippers2022}. The PP application maps each strength to its corresponding
impairment and designs an interaction module that activates the strength to
naturally redirect attention toward the required task.

\section{Related Works}

The neurodiversity of executive functions in ADHD is concentrated in four brain
regions: the prefrontal cortex, the basal ganglia, the cerebellum, and the
parietal-temporal region, operating through the reticular activating
system~\cite{castellanos2002,valera2007}. These regions cannot be categorised
independently; they form a dynamic, integrated management process in which
executive processing depends on coordinated activity across multiple cognitive
components. This section analyses that neurodiversity by (i)~examining the
neural connections through which positive and negative characteristics can be
interchanged, (ii)~reviewing functional modules of existing ADHD management
applications, and (iii)~reviewing the interaction design principles required for
mobile application development. All three inform the PP prototype.

\subsection{Research Overview}

A fundamental assumption of the neurodiversity paradigm is that all forms of
neurological variation are valuable~\cite{denhouting2019}. Within this paradigm,
qualitative research by Sedgwick et al.~\cite{sedgwick2019} and Schippers
et al.~\cite{schippers2022} identifies six core positive themes in adults with
ADHD: \textit{cognitive vitality}, \textit{courage}, \textit{energy},
\textit{humanity}, \textit{resilience}, and \textit{transcendence}. Brown's
executive functioning model~\cite{nigg2005} identifies six corresponding
impairment dimensions---activation, concentration, effort, emotion, memory, and
execution. The six dimensions within this model map one-to-one
onto the six positive themes, a structural correspondence arising from shared
neural substrates: the same brain regions and networks that produce each
executive impairment also generate the corresponding cognitive strength. Each
strength is therefore a principled, neurologically grounded compensatory tool
for its paired deficit. The six pairings and their neural bases are as follows.

\textbf{Cognitive Vitality vs.\ Concentration.}
Cognitive vitality summarises the uninterrupted mental activity, spontaneous
non-sequential thought, and intense internal concentration characteristic of
ADHD. The corresponding neural system is the brain's default mode network
(DMN), which is negatively activated during external attention tasks and
positively activated during internal tasks related to self, memory, and future
vision~\cite{sestieri2011}. In ADHD the DMN remains persistently active,
producing the inability to focus on external tasks for long periods while
simultaneously preventing disengagement from internal ideation. \textit{Design implication:} Divergent thinking, creativity, and curiosity may provide alternative entry points for redirecting attention from internally generated thoughts towards externally required tasks..

\textbf{Courage vs.\ Emotion.}
Courage summarises how people with ADHD face fears and deal with uncertainty
through non-conformity, risk-taking, and the drive for individualisation. The
neural system involved is the amygdala (over-activated in ADHD, including
dopamine receptor density) in conjunction with the ventral striatum
(under-activated)~\cite{knutson2001}. This produces emotional impulsivity,
inadequate self-regulation, and aversive responses to long-term reward
anticipation. \textit{Design implication:} Sensitivity to uncertainty, novelty, and immediate rewards suggests that task engagement may be supported through positively framed, short-horizon actions rather than distant or aversive outcomes..

\textbf{Energy vs.\ Activation.}
Energy summarises the psycho-energetic willpower motivated by a strong desire to
achieve something meaningful. The dorsal attention network (DAN) is more active
than the ventral attention network (VAN) in ADHD, elevating top-down selective
attention while impeding bottom-up priority shifting~\cite{corbetta2002}.
Consequently, ADHD does not perceive time as a sequence but as a dispersed
collection of emotionally linked events, making task prioritisation difficult.
Yet hyper-focus emerges when an ambitious personal goal is present.
\textit{Design implication:} Because task engagement may increase when activities are connected to personally meaningful goals, design should emphasise purpose and personal significance rather than relying solely on conventional priority ordering.

\textbf{Humanity vs.\ Execution.}
Humanity summarises social intelligence, humour, self-acceptance, and empathy
in social behaviour. Mirror neurons produce strong empathic resonance and
positive emotional contagion but also difficulty adapting behaviour when
others are puzzled or upset. \textit{Design implication:} Empathy, humour, and social connectedness may be used to support positive affect and sustained engagement, although social motivation alone may not be sufficient to guide task execution..

\textbf{Resilience vs.\ Effort.}
Resilience describes how people with ADHD self-regulate and self-protect under
adverse conditions through behavioural adjustment strategies. The anterior
cingulate cortex drives cognitive control and behavioural monitoring, but ego
depletion under antagonistic conflict leads to mental fatigue and ineffective
sustained self-regulation~\cite{carter2007}. \textit{Design implication:}
Adaptive coping strategies should be supported as flexible and optional resources that users can draw upon during periods of cognitive fatigue, rather than as fixed or compulsory procedures.

\textbf{Transcendence vs.\ Memory.}
Transcendence describes the rapid aesthetic absorption---awe, wonder, and deep
memorisation---that people with ADHD experience, particularly with music and
art. The corticocerebellar network with norepinephrine underpins working memory;
in ADHD, internal norepinephrine links are weakened, yet high absorption during
aesthetic experiences produces deep encoding~\cite{shiota2006}.
\textit{Design implication:} Aesthetic engagement and deep absorption may support the encoding and retention of task-related information, suggesting value in visually or emotionally meaningful forms of representation.

\subsection{Mobile Tools Research}

\subsubsection{User Interaction Design.}
Mental models are the key to user-centred design~\cite{ruiz2021}. A mental
model is a person's abstract concept of a device that enables them to interpret
and predict an application. Research shows that graphical interfaces outperform
text-only representations, and that visual, symbolic, and pictorial design
elements are the primary medium through which mental models are
conceptualised. From a systematic review of 257 widely cited UI design
principles, six core guidelines were identified: providing informative feedback,
minimising memory load, clear error messages, providing shortcuts, using natural
language, and using real-world metaphors~\cite{ruiz2021}. The inherent mental
model of people with ADHD is largely negative (``I need to be managed''); PP's
framework must extend this model with positive elements through a virtuous-circle
design logic.

\subsubsection{Interactive Modules.}
A review of 22 representative ADHD management applications (10,000--50,000
downloads) was conducted, covering diagnostic tools, task management, habit
building, time management, focus enhancement, emotion regulation, memory
support, and financial management~\cite{pasarelu2020}. Table~\ref{tab:apps}
summarises the resulting module analysis by dysfunction. The review reveals that
most modules address symptoms through mandatory reminders and timers with no
leverage of ADHD cognitive strengths, while a subset---mapping, meditation,
community interaction, and random creative prompts---do align with positive ADHD
characteristics. No existing tool integrates these strength-aligned modules into
a unified framework.

\begin{table*}
  \caption{ADHD application module analysis: dysfunction targeted, interaction
    type, and whether the module leverages ADHD cognitive strengths (adapted
    from review of 22 representative apps).}
  \label{tab:apps}
  \begin{tabular}{p{0.10\linewidth} p{0.16\linewidth} p{0.20\linewidth}
                  p{0.30\linewidth} c}
    \toprule
    \textbf{Dysfunction} & \textbf{Symptom} & \textbf{Module} &
    \textbf{Interaction Element} & \textbf{Leverages Strengths?} \\
    \midrule
    Activation  & Time chaos            & Task reminder    & Mandatory reminder          & No  \\
    Activation  & Time chaos            & Label / tagging  & Information capture          & Yes \\
    Focus       & Lack of attention     & Timer            & Mandatory countdown          & No  \\
    Effort      & Inability to sustain  & Timer            & Mandatory countdown          & No  \\
    Emotion     & Impulsive decision    & Meditation       & Mindful relaxation           & Yes \\
    Memory      & Missing tasks         & Mapping          & Visual spatialisation        & Yes \\
    Execution   & Cannot adjust behaviour & Habit tracker  & Self-awareness, personalised time & Yes \\
    Social      & Low positive mood     & Community        & Patient interaction          & Yes \\
    Creativity  & Dispersed thinking    & Random prompts   & Idea generation              & Yes \\
    \bottomrule
  \end{tabular}
\end{table*}

\subsection{Research Conclusions}

The background analysis yields eight design conclusions that inform PP:

\begin{enumerate}
  \item Use divergent thinking and creativity to redirect ADHD attention from
    intrinsic tasks toward the required extrinsic task.
  \item Replace long-term reward tasks with chains of short-term positive
    rewards to evoke impulsive decision-making without triggering anxiety.
  \item Replace abstract priority ordering with a high-meaning goal narrative;
    completing the priority task becomes the necessary path to a personally
    significant outcome.
  \item Present memory-intensive tasks in an artistic form to reduce working
    memory loss.
  \item Add a community section to sustain positive mood through social
    connection beyond the individual patient.
  \item Draw on strength-aligned modules identified in the review: tagging,
    positive thinking, mapping, habit tracking, community interaction, and
    random creative prompts.
  \item Extend the inherent negative mental model (``I need to be managed'')
    with positive elements, using a virtuous circle to motivate self-initiated
    action.
  \item Adhere to the six core UI principles and use visual, symbolic, and
    pictorial design to conceptualise a positive mental model.
\end{enumerate}

\section{Theoretical Framework: Six Strengths and Six Deficits}

The following pairings form the conceptual core of PP.

\textbf{Cognitive Vitality vs.\ Concentration.} The brain's default network is
persistently active in ADHD, causing mind-wandering yet enabling divergent
thinking and creativity. The neural system involved produces spontaneous,
non-sequential ideation. Divergent creativity is therefore used as a hook to
redirect attention toward external tasks.

\textbf{Courage vs.\ Emotion.} The amygdala is over-activated in ADHD, leading
to emotional impulsivity but also risk-taking and non-conformity. The
over-sensitivity to negative stimuli produces procrastination. Positive
risk-taking courage is leveraged to displace avoidant behaviour.

\textbf{Energy vs.\ Activation.} The dorsal attention network (DAN) dominates in
ADHD, generating hyper-focus when a personally meaningful goal is present but
causing neglect of lower-priority tasks~\cite{corbetta2002}. High-meaning
goal-setting replaces traditional priority ordering to activate task engagement.

\textbf{Humanity vs.\ Execution.} Mirror neuron activity drives strong social
empathy and humour. Positive social connections sustain emotional regulation,
extending this strength to maintain a supportive mood during task execution.

\textbf{Resilience vs.\ Effort.} The anterior cingulate cortex drives
self-regulation, but ego depletion leads to mental fatigue and ineffective
sustained effort. Resilience is channelled into adaptive coping strategies that
are embedded as options within the application flow.

\textbf{Transcendence vs.\ Memory.} People with ADHD show high absorption
during aesthetic experiences, facilitating deep memory encoding. Tasks that
require working memory are therefore presented in an artistic, visually rich
form to strengthen retention.

Table~\ref{tab:framework} summarises the six pairings together with the
responsible neural system and the resulting design direction for each module.

\begin{table*}
  \caption{Neurodiversity analysis: cognitive strengths, executive deficits,
    neural systems, and PP design directions.}
  \label{tab:framework}
  \begin{tabular}{p{0.10\linewidth} p{0.11\linewidth} p{0.18\linewidth} p{0.50\linewidth}}
    \toprule
    \textbf{Strength} & \textbf{Executive Deficit} & \textbf{Neural System} & \textbf{Design Direction} \\
    \midrule
    Cognitive Vitality &
      Concentration &
      Default Mode Network (persistently active; internal ideation) &
      Harness divergent thinking and creativity as an entry point for external tasks; a creative word-association exercise focuses attention on the required task subject. \\
    \addlinespace
    Courage &
      Emotion &
      Amygdala (over-activated) and Ventral Striatum (under-activated) &
      Leverage non-conformity, risk-taking, and persistence to displace procrastination; replace aversive long-term rewards with chains of short-term positive rewards. \\
    \addlinespace
    Energy &
      Activation &
      Dorsal Attention Network (DAN dominates over VAN) &
      Replace abstract priority ordering with high-meaning goal narratives; associating a task with a personally significant purpose activates willpower to begin. \\
    \addlinespace
    Humanity &
      Execution &
      Mirror Neurons (strong empathic resonance) &
      Use positive social connection as an extension tool; community sharing sustains positive mood and supports behaviour continuity. \\
    \addlinespace
    Resilience &
      Effort &
      Anterior Cingulate Cortex (self-regulation limited by ego depletion) &
      Embed adaptive coping strategies as optional branching choices rather than mandatory steps; resilience supports choice-driven self-management. \\
    \addlinespace
    Transcendence &
      Memory &
      Corticocerebellar Network with Norepinephrine (deep aesthetic encoding) &
      Present memory-intensive task information in an artistic, puzzle-based form; aesthetic engagement deepens working memory encoding. \\
    \bottomrule
  \end{tabular}
\end{table*}

A review of 22 representative ADHD management applications (10,000--50,000
downloads) confirmed that no existing tool exploits these positive
characteristics; all rely on mandatory reminders, timers, or correction
mechanisms~\cite{pasarelu2020}.

\section{User Group and Requirements}

Approximately 2.8\% of adults worldwide have ADHD~\cite{faraone2006}. An
analysis of two qualitative studies (422 adults with ADHD across six countries)
showed that participants described their own strengths---problem-solving, fantasy,
novelty-seeking, deep focus, empathy, resilience---yet received little external
recognition of them. A review of 127 ADHD studies found that 29\% of treated and
56\% of untreated adults reported low self-esteem, and 38--73\% felt their social
functioning was affected~\cite{harpin2016}.

PP therefore targets two parallel needs: (1)~social recognition and affirmation
of ADHD strengths; and (2)~support for executive skill and self-esteem through
positive engagement rather than coercion.

\section{Design}

\subsection{Module Logic}

Each of the six cognitive strengths maps to a distinct application module:

\begin{itemize}
  \item \textbf{Concentration module} (Cognitive Vitality): The user flips five
    favourite images and clicks to reveal five random words. Their task is to
    create something based on those words with the required external task as the
    subject, harnessing divergent thinking to pull attention toward the task.

  \item \textbf{Memory module} (Transcendence): Before a puzzle task, the user
    enters information they need to retain. Each item is mapped to a puzzle
    piece; hovering over a piece in the puzzle displays the corresponding fact,
    reinforcing working memory through aesthetic engagement.

  \item \textbf{Persistent tasks module} (Humanity): The user occupies the centre
    of a solar-system visualisation; each ongoing habit or responsibility is a
    planet that lights up upon completion, building a positive social-identity
    narrative.

  \item \textbf{Long-term tasks module} (Energy): Long-term tasks are broken into
    short-term sub-tasks, each linked to an immediate reward, converting aversive
    distant goals into a chain of motivating short-horizon achievements.

  \item \textbf{Priority tasks module} (Courage): Users simulate costumes and
    scenes from a favourite film or TV production to role-play completing
    priority tasks in reality, anchoring task execution in a high-meaning,
    self-determined narrative.

  \item \textbf{Emotion module} (Resilience): When users are in a positive mood,
    social expression links let them help other users or leave feedback; when
    in a negative mood, a positive-thinking and meditation module is available.
\end{itemize}

\subsection{Workflow}

The application workflow follows a non-linear model: users enter whichever
module matches their current cognitive state rather than following a prescribed
sequence. This respects the temporal experience of ADHD, in which time is
perceived as a dispersed collection of events linked to people and emotions
rather than a sequential schedule.

\subsection{Visual Design Principles}

People with ADHD have heightened visuospatial processing sensitivity. The PP
interface therefore uses: non-repetitive, non-patterned layouts; vivid colours
representing high energy and positivity; irregular icon shapes that echo
creativity and divergent thinking; and unique artwork in puzzle sections. The
UI adheres to six empirically validated core principles: informative feedback,
minimal memory load, good error messages, shortcuts, natural language, and
real-world metaphors~\cite{ruiz2021}.

\section{Prototype}

A high-fidelity prototype was developed from a wireframe framework that
implements the module logic described above. The visual system consistently
applies the design principles across all modules, ensuring that aesthetic
stimulation supports---rather than distracts from---each cognitive strength
activation. The prototype is interactive and supports the full task-entry-to-
completion flow for each module.

\section{User Testing}

Preliminary evaluation was conducted with three anonymous adults with ADHD.
Positive feedback: the application is engaging and the colour palette
successfully attracts and holds attention. Critical finding: in noisy external
environments (e.g.\ commuting), ambient auditory distractions competed with the
visual engagement. As a result, the next iteration of PP will integrate
background music and sound effects. Prior research shows that neurofeedback audio
protocols are a well-established adjunctive treatment for ADHD~\cite{enriquez2019},
providing a principled basis for this extension.

\section{Conclusion}

PP is, to our knowledge, the first mobile application designed around ADHD
cognitive strengths to circumvent executive disadvantages. It replaces coercion,
reminders, and correction with creativity, high energy, dispersed thinking, and
high-output engagement. The neurodiversity paradigm underlying PP provides a
replicable design framework: identify the neural pathway linking a positive ADHD
characteristic to its corresponding executive impairment, then design an
interaction that activates the strength while naturally redirecting the
impairment.

Limitations include a small user testing sample (N\,=\,3) and limited coverage of
users with deafblindness. Future work will expand evaluation and incorporate
neurofeedback technology as wireless EEG device costs continue to fall. Beyond
task management, PP aims to shift the mental model of ADHD from ``I need to be
managed'' to one that affirms capability and positive identity.

\begin{acks}
The author thanks the three anonymous ADHD participants who provided user testing
feedback. This work was conducted under UCL ethical guidelines; participants
submitted informed consent via the UCL REDCap platform.
\end{acks}

\section*{Ethics and Privacy Statement}

User testing involved three anonymous adults with ADHD. No personally
identifiable information was collected. Participants provided informed consent
through the UCL REDCap platform in accordance with UCL ethical review guidelines.
The study presents minimal societal risk; the design aims to benefit ADHD users
by affirming their strengths rather than pathologising their differences.

\bibliographystyle{ACM-Reference-Format}
\bibliography{adhd-poster}

\end{document}
\endinput
